\documentclass[11pt]{article}
\usepackage[margin=0.65in]{geometry}
\usepackage{booktabs}
\usepackage{float}
\usepackage[T1]{fontenc}
\usepackage[utf8]{inputenc}
\begin{document}
\section*{Input-output price shocks and inflation inequality}
\begin{table}[H]
\centering
\caption{Effect of a 40\% Energy-Price Shock with Cross-Country Propagation}
\scriptsize
\setlength{\tabcolsep}{1.2pt}
\begin{tabular}{lrrrrrrr}
\toprule
Country & \shortstack{Effect on average\\inflation -- Total} & \shortstack{Direct effect on\\average inflation} & \shortstack{Indirect effect on\\average inflation} & \shortstack{Indirect effect on average\\inflation through trade} & \shortstack{Effect on\\inflation gap} & \shortstack{Direct effect on\\inflation gap} & \shortstack{Indirect effect on\\inflation gap}
 \\
\midrule
Austria & 3.3 & 1.9 & 1.4 & 0.8 & 0.7 & 0.5 & 0.2
 \\
Belgium & 4.6 & 2.6 & 1.9 & 1.2 & 1.9 & 1.7 & 0.2
 \\
Cyprus & 4.0 & 3.3 & 0.7 & 0.7 & 2.0 & 1.7 & 0.3
 \\
Germany & 4.1 & 3.0 & 1.1 & 0.3 & 1.5 & 1.4 & 0.1
 \\
Estonia & 4.4 & 2.1 & 2.3 & 0.9 & 2.2 & 0.8 & 1.4
 \\
Greece & 3.6 & 2.8 & 0.9 & 0.2 & 2.1 & 1.9 & 0.2
 \\
Spain & 3.1 & 2.2 & 0.9 & 0.3 & 0.2 & 0.1 & 0.1
 \\
Finland & 2.9 & 2.0 & 0.8 & 0.2 & -1.1 & -1.0 & -0.1
 \\
France & 3.2 & 2.3 & 0.8 & 0.3 & 0.6 & 0.6 & -0.0
 \\
Croatia & 6.0 & 4.3 & 1.6 & 0.8 & 5.0 & 3.4 & 1.6
 \\
Ireland & 2.2 & 1.8 & 0.4 & 0.1 & 1.8 & 1.6 & 0.1
 \\
Italy & 5.1 & 3.4 & 1.7 & 0.2 & 1.7 & 1.5 & 0.2
 \\
Lithuania & 4.1 & 3.5 & 0.7 & 0.3 & 0.8 & 0.9 & -0.0
 \\
Luxembourg & 5.0 & 4.6 & 0.4 & 0.4 & 0.4 & 0.4 & 0.0
 \\
Latvia & 4.3 & 3.1 & 1.2 & 1.2 & 3.3 & 2.3 & 1.0
 \\
Malta & 2.2 & 1.7 & 0.5 & 0.5 & 0.1 & 0.1 & -0.0
 \\
Netherlands & 3.7 & 2.9 & 0.8 & 0.4 & 1.4 & 1.4 & -0.0
 \\
Portugal & 3.0 & 2.0 & 1.1 & 0.5 & 1.1 & 0.9 & 0.2
 \\
Slovenia & 4.5 & 3.9 & 0.6 & 0.5 & 1.8 & 1.7 & 0.1
 \\
Slovakia & 2.4 & 1.7 & 0.7 & 0.3 & 0.5 & 0.5 & 0.0
 \\
\midrule
\textbf{Euro area} & \textbf{3.9} & \textbf{2.7} & \textbf{1.1} & \textbf{0.4} & \textbf{1.2} & \textbf{1.1} & \textbf{0.1}
 \\
\bottomrule
\end{tabular}
\par\vspace{0.35em}\footnotesize\noindent Sources: Eurostat HICP, HBS, FIGARO, authors' calculations.
\par\vspace{0.25em}\footnotesize\noindent Notes. Entries are percentage-point effects of a uniform 40 percent energy-price shock. The price response is $\Delta p = (I - A^{\prime})^{-1}s$, where $A_{(o,i),(d,j)}=Z_{(o,i),(d,j)}/x_{d,j}$ retains bilateral intermediate flows from origin country $o$ and sector $i$ to destination country $d$ and sector $j$. Direct effects arise from household consumption weights mapped to the shocked inputs; indirect effects include domestic and cross-country propagation through intermediate imports and exports. The trade column is the difference between the full EA20 solution and a counterfactual solution in which all off-diagonal country blocks of $A$ are set to zero; it is included in, rather than added to, the total indirect effect. The euro-area aggregate covers all 20 member countries. Imports from origins outside the EA20 production system enter production costs but are treated as exogenous and do not generate additional endogenous feedback. FIGARO flows are measured in current million euros at basic prices; the dataset used here does not provide separate import or export deflators.
\end{table}
\vspace{0.6em}
\end{document}
